\chapter{Introducción}
\label{ch:intro}

El caso asignado al Grupo~1 es la secuencia R2ET. Un robot industrial reparte de forma alternada, entre dos esteras transportadoras, las piezas que salen de una máquina de proceso; ninguna estera admite más de tres piezas a la vez. El control resuelve dos problemas al tiempo: mantener el llenado de cada estera cerca del punto de operación ($1.5$~piezas) y repartir el robot para que ninguna se vacíe ni se sature, incluso cuando cambia el peso de las piezas.

Sobre esa misma planta y los mismos escenarios de perturbación se comparan cuatro estrategias. Un PID discreto sirve de referencia industrial. Frente a él se ponen un controlador difuso de 35 reglas lingüísticas auditables, una red neuronal profunda (\emph{Deep Learning}, DL) entrenada \emph{offline} mediante retropropagación a través de la simulación discreta de la planta, y un agente de aprendizaje por refuerzo tabular (\emph{Q-Learning}, QL) que aprende por interacción con el entorno.

La evaluación se organiza en dos fases complementarias. La Fase~1 prueba cada método como servorregulador sobre el modelo reducido de primer orden obtenido de la dinámica dominante de la planta asignada al Grupo~1, con escalón de referencia y perturbación de carga acotada; así se miden las prestaciones servo al margen de la complejidad del sistema real. La Fase~2 evalúa el sistema R2ET completo: dos esteras, robot compartido, pulso de alta demanda y atasco parcial por cambio de peso. En esta fase se mide la coordinación supervisora.

\footnote{Se emplea el punto como separador decimal por coherencia con el banco de simulación y las figuras generadas.}

Todas las simulaciones corren en tiempo discreto, con parámetros fijos y reproducibles. El Apéndice~A incluye extractos representativos de las funciones MATLAB y de la lógica empleada para generar los modelos, las capturas Simulink y las figuras de resultados.
